\documentclass[10pt,a4paper]{article}
\usepackage[T1]{fontenc}
\usepackage{lmodern}
\usepackage[spanish]{babel}
\usepackage[margin=1.6cm]{geometry}
\usepackage{microtype,booktabs,tabularx,xcolor}
\usepackage[colorlinks=true,urlcolor=blue!50!black]{hyperref}
\setlength{\parindent}{0pt}\setlength{\parskip}{0.3em}
\pagestyle{empty}
\newcolumntype{Y}{>{\raggedright\arraybackslash}X}
\begin{document}
{\large\textbf{Calor y sequ\'ia en Catalu\~na: resumen de una p\'agina}}\hfill{\small Gabriel Albasini Garaulet, noviembre de 2026}

\smallskip
\textbf{Idea.} Tres preguntas sencillas sobre efectos econ\'omicos y sociales del calor y la sequ\'ia, respondidas solo con datos abiertos. Empezamos con cuatro expectativas: tres fallaron y una se cumpli\'o d\'ebilmente.

\smallskip
\small
\begin{tabularx}{\linewidth}{@{}p{2.1cm}YY@{}}
\toprule
\textbf{Tema} & \textbf{Qu\'e encontramos} & \textbf{L\'imites} \\
\midrule
\textbf{Demanda el\'ectrica y temperatura} (idea 6)
& Zona de confort de 17 a 22\,$^\circ$C (\(\approx\)114 GWh/d\'ia). Cada grado por encima suma \textbf{4,2 GWh} (3,7\,\%); cada grado por debajo, \textbf{2,9 GWh}. El calor pesa m\'as que el fr\'io (esper\'abamos lo contrario). En 912 d\'ias laborables de 2022--2025.
& Una estaci\'on para toda Catalu\~na. Con Fabra (Collserola, 2,3\,$^\circ$C m\'as fr\'ia) los umbrales bajan a 15/19\,$^\circ$C y la pendiente de calor a 3,2. La temperatura explica una parte, no todo (error 5,8 GWh). \\
\addlinespace
\textbf{Sequ\'ia e hidroel\'ectrica} (idea 6)
& Los embalses cayeron del 95\,\% al 15\,\% (feb. 2024), pero \textbf{el ciclo combinado no cubri\'o el hueco}: en 2023 baj\'o tambi\'en (\(-2\,014\) GWh) y las \textbf{importaciones netas subieron 2\,575 GWh} (de 3\,098 a 5\,672, el 13\,\% de la demanda). La hidr\'aulica sigue a la \textbf{lluvia del Pirineo} (r = 0,80 anual) m\'as que a los embalses usados (0,69), que son de otra cuenca.
& La hidr\'aulica es solo el 5,8\,\% de la generaci\'on. Correlaciones, no causas. Endesa atribuye su +60\,\% de 2024 a m\'as agua (+20\,\%) y a mejor operaci\'on. \\
\addlinespace
\textbf{Aceite de oliva y sequ\'ia} (idea 4)
& El aceite subi\'o un \textbf{54,6\,\%} en dic. 2023 (m\'aximo +73,5\,\% en oct.). Con la lluvia de Andaluc\'ia: r = \(-0{,}42\) (p = 0,047, 23 a\~nos); con la catalana: \(-0{,}31\) (p = 0,15). Tres de los cuatro a\~nos m\'as secos de Andaluc\'ia vieron subidas de precio.
& \textbf{Evidencia d\'ebil}: significativa solo con Pearson y por poco; con Spearman no. Otras subidas (2015, 2021--22) ocurrieron con sequ\'ias suaves. La lluvia andaluza es de rean\'alisis, no de estaciones. \\
\addlinespace
\textbf{Refugios clim\'aticos en Barcelona} (idea 2)
& El \textbf{99,5\,\%} de los residentes est\'a a menos de 800 m de un refugio (96,6\,\% a 400 m); igual para mayores de 75. Los barrios pobres \textbf{no} est\'an peor servidos: el decil m\'as pobre tiene el 97\,\% a 400 m. El m\'as rico tiene la menor cobertura (88,3\,\%) por su baja densidad. Los refugios por 10\,000 habitantes son casi iguales en secciones de vulnerabilidad baja, media y alta (3,1/3,4/3,3).
& Distancia en l\'inea recta (la ruta real es m\'as larga). Los refugios son muy distintos (piscinas, bibliotecas, patios\ldots). La vulnerabilidad es del estudio de 2015, con poca documentaci\'on. Acceso no es exposici\'on: no medimos temperatura. \\
\bottomrule
\end{tabularx}

\normalsize
\smallskip
\textbf{Qu\'e decir ma\~nana (tres frases).}
\begin{itemize}\setlength{\itemsep}{0pt}\setlength{\topsep}{2pt}
\item El calor sube m\'as la demanda el\'ectrica que el fr\'io, y el resultado cambia un poco seg\'un el term\'ometro que se use.
\item La sequ\'ia no la cubri\'o el gas: la cubrieron las importaciones, y los embalses cercanos a Barcelona no son los que mueven la hidroel\'ectrica.
\item En Barcelona el acceso a un refugio no depende de la renta; la pregunta pendiente es la exposici\'on real al calor.
\end{itemize}
\textbf{Pr\'oximos pasos.} Distancias a pie por la red de calles; temperatura medida en la ciudad; datos de cosecha de aceite; embalses del Pirineo (no encontrados en datos abiertos catalanes).

{\footnotesize\raggedright Todo el c\'odigo y las figuras est\'an en \url{https://github.com/NuriaMayoral/XEMAS-} (\texttt{figures/Gabri}). Art\'iculo completo en \texttt{paper/gabri\_paper.pdf}.\par}
\end{document}
